\section{FOREIGN KEY dan Referential Integrity}

\textbf{Foreign key} (FK) adalah constraint yang memastikan nilai di satu kolom merujuk ke nilai yang ada di primary key tabel lain. Ini menjaga \textbf{referential integrity} --- menjamin bahwa tidak ada data ``yatim'' (orphaned records) yang merujuk ke data yang tidak ada. Misalnya, FK pada kolom \code{kategori\_id} di tabel \code{produk} memastikan bahwa setiap produk merujuk ke kategori yang benar-benar ada di tabel \code{kategori} \cite{mysqlDocs}.

MySQL menyediakan aksi referensial yang menentukan apa yang terjadi ketika baris yang dirujuk dihapus atau diubah: \code{CASCADE} (ikut dihapus/diubah), \code{SET NULL} (diset NULL), \code{RESTRICT} atau \code{NO ACTION} (tolak operasi --- default), dan \code{SET DEFAULT}. Pemilihan aksi yang tepat sangat penting: \code{CASCADE} praktis untuk data yang benar-benar bergantung (misalnya detail pesanan saat pesanan dihapus), sedangkan \code{RESTRICT} lebih aman untuk mencegah penghapusan yang tidak disengaja \cite{mysqlGettingStarted}.

\begin{webcode}[language=SQL, caption={FOREIGN KEY dengan aksi referensial}]
CREATE TABLE pesanan (
  id      INT AUTO_INCREMENT PRIMARY KEY,
  user_id INT NOT NULL,
  tanggal DATETIME DEFAULT CURRENT_TIMESTAMP,
  total   DECIMAL(12,2) DEFAULT 0,
  FOREIGN KEY (user_id) REFERENCES pengguna(id)
    ON DELETE RESTRICT
    ON UPDATE CASCADE
);

CREATE TABLE pesanan_detail (
  id         INT AUTO_INCREMENT PRIMARY KEY,
  pesanan_id INT NOT NULL,
  produk_id  INT NOT NULL,
  jumlah     INT NOT NULL DEFAULT 1,
  harga      DECIMAL(12,2) NOT NULL,
  FOREIGN KEY (pesanan_id) REFERENCES pesanan(id) ON DELETE CASCADE,
  FOREIGN KEY (produk_id) REFERENCES produk(id) ON DELETE RESTRICT
);
\end{webcode}

\begin{catatan}
Foreign key hanya berfungsi pada engine \textbf{InnoDB} (default di MySQL modern). Engine MyISAM tidak mendukung foreign key. Pastikan tabel menggunakan InnoDB dengan menambahkan \code{ENGINE=InnoDB} saat CREATE TABLE atau set default engine di konfigurasi server.
\end{catatan}

